% (C) 2026 Ralf Stephan, in collaboration with Claude Code.  CC0 1.0.
\chapter{Splittings, homomorphisms and products}\label{ch:splitting}

The same problem can be encoded by many carry graphs. The \texttt{DubC} engine works modulo
\(M=14M'\) or modulo \(M'\). The interval problem uses residues, cells, or both at once. Word
graphs with long edges can be broken into letters. This chapter collects the constructions that
change the graph and records what they preserve. The invariants that matter are the entropy,
zero entropy (Theorem~\ref{thm:six-faces}), an empty core (Theorem~\ref{thm:core-empty}) and
the label form (Definition~\ref{def:allper}). The vertex count is never one of them.

\section{Right covers}

\begin{definition}\label{def:rightcover}
Let \(X\) and \(Y\) be carry graphs on \(V\) and \(W\), and \(U\subseteq W\). A map
\(p\colon W\to V\) is a \emph{right cover}\index{right cover} of \(X\) by \(Y\) on \(U\) if \(U\) is closed under
successors and, for \(y\in U\) and every label \(a\), \(p\) maps the \(a\)-successors of \(y\)
bijectively onto the \(a\)-successors of \(p(y)\).
\end{definition}

This is the graph form of a right-resolving factor code \cite[§8.2]{LM95}. A right cover on all
of \(W\) is a homomorphism \(Y\to X\).

\begin{proposition}\label{prop:rightcover-paths}
If \(p\) is a right cover on \(U\), then for \(y\in U\) and all \(n\)
\[
\mathrm{npaths}_Y(y,n)=\mathrm{npaths}_X(p(y),n).
\]
\end{proposition}

\begin{proof}
Induction on \(n\). The recursion
\[
\mathrm{npaths}(y,n+1)=\sum_a\sum_{z\in\succc(y,a)}\mathrm{npaths}(z,n)
\]
is carried term by term along the bijections.
\end{proof}

A vertex is a \emph{sink} if no arc leaves it.

\begin{theorem}\label{thm:rightcover}\index{right cover}
Let \(V\), \(W\) and \(A\) be finite and \(p\) a right cover on \(U\). Suppose every vertex of
\(W\setminus U\) is a sink, and every vertex of \(V\) is a sink or has a preimage in \(U\). Then
\(X\) and \(Y\) have the same entropy, both or neither have zero entropy, and both or neither
have an empty core.
\end{theorem}

\begin{proof}
For \(n\ge1\), sinks carry no path of length \(n\). Choosing one preimage of each non-sink vertex
of \(X\) and using Proposition~\ref{prop:rightcover-paths} gives
\(\mathrm{ntot}_X(n)\le\mathrm{ntot}_Y(n)\). Bounding each \(\mathrm{npaths}_Y(y,n)\) by
\(\mathrm{ntot}_X(n)\) gives \(\mathrm{ntot}_Y(n)\le|W|\,\mathrm{ntot}_X(n)\). Total path counts
comparable up to constant factors have the same exponential growth rate, and they vanish at
the same lengths. Then apply Theorem~\ref{thm:six-faces} (entropy \(\le0\)) and
Theorem~\ref{thm:core-empty}.
\end{proof}

\begin{example}[full and reduced model]\label{ex:full-reduced}
For base \(7\) the digit graph \(D_{7,14M'}\) of Example~\ref{ex:digit} has the units modulo
\(14M'\) as vertices. Its \emph{reduced model} \(R_{M'}\) has the units modulo \(M'\) and arcs
\(\rho\xrightarrow{d}7\rho+d\) for the digits \(d\in\{2,4,6\}\). These are the digits that keep
\(7r+d\) prime to \(2\) and \(7\). Reduction modulo \(M'\) is a right cover on the units. The
\(d\)-successor \(7x+d\) of a unit \(x\) is a unit iff \(d\in\{2,4,6\}\) and \(7x+d\) is prime to
\(M'\), and this is exactly the condition for the \(d\)-arc of \(R_{M'}\) at the reduction. Each
side has at most one \(d\)-successor, so the bijection is automatic. If \(\gcd(M',14)=1\), every
unit modulo \(M'\) lifts to a unit modulo \(14M'\) by the Chinese remainder theorem. Theorem~\ref{thm:rightcover}
then says that the two models have the same zero-entropy verdict. Their cores still differ in
size: at \(y=19\) they have \(53\,218\) and \(48\,548\) vertices. The full core is an in-splitting
of the reduced one.
\end{example}

\section{State splitting}

Let \(\pi\colon V\times A\times V\to C\) assign a class to every arc, with \(C\) finite and
non-empty.

\begin{definition}\label{def:split}
The \emph{in-splitting}\index{in-splitting} \(X_\pi^{\mathrm{in}}\) \cite[Def.~2.4.7]{LM95} has vertices \((u,c)\)
and an arc \((u,c)\xrightarrow{a}(v,\pi(u,a,v))\) for every arc \(u\xrightarrow{a}v\) of \(X\):
the copies of \(v\) sort the arcs \emph{into} \(v\). The \emph{out-splitting}\index{out-splitting}
\(X_\pi^{\mathrm{out}}\) \cite[Def.~2.4.3]{LM95} has an arc \((u,c)\xrightarrow{a}(v,c')\) for
every arc \(u\xrightarrow{a}v\) of class \(\pi(u,a,v)=c\) and every \(c'\): the copies of \(u\)
share out the arcs \emph{out of} \(u\).
\end{definition}

\begin{theorem}\label{thm:split}
\begin{enumerate}[label=(\alph*)]
\item The first projection is a right cover of \(X\) by \(X_\pi^{\mathrm{in}}\) on all
  vertices. So \(\mathrm{npaths}(\,(u,c),n)=\mathrm{npaths}(u,n)\).
\item \(\sum_{c\in C}\mathrm{npaths}_{X_\pi^{\mathrm{out}}}((u,c),n)=|C|\cdot
  \mathrm{npaths}_X(u,n)\).
\item Both splittings have the entropy of \(X\), have zero entropy iff \(X\) has, and have an
  empty core iff \(X\) has.
\end{enumerate}
\end{theorem}

\begin{proof}
(a) The \(a\)-successors of \((u,c)\) are the pairs \((v,\pi(u,a,v))\), one for each
\(a\)-successor \(v\) of \(u\). (b) Induction on \(n\). A path from \((u,c)\) starts with an arc
of class \(c\), and summing over \(c\) recovers every arc out of \(u\) exactly once. (c) From
(a) and Theorem~\ref{thm:rightcover} for the in-splitting. From (b),
\(\mathrm{ntot}_{X_\pi^{\mathrm{out}}}=|C|\,\mathrm{ntot}_X\) for the out-splitting.
\end{proof}

Part (b) is the counting content of the conjugacy theorem for splittings
\cite[Thm.~2.4.10]{LM95}, which needs no conjugacy here. A right cover need not be surjective
or injective on vertices. The vertex count is free while every growth invariant is fixed.

\section{Line graphs and higher blocks}

\begin{definition}\label{def:linegraph}
The \emph{line graph}\index{line graph} \(X^{[2]}\) \cite[§4.5]{BJG09}, the edge graph of
\cite[Def.~2.3.4]{LM95}, has the triples \((u,a,v)\) as vertices and an arc
\((u,a,v)\xrightarrow{b}(v,b,w)\) whenever both \(u\xrightarrow{a}v\) and \(v\xrightarrow{b}w\)
are arcs of \(X\). Triples that are not arcs are isolated. Iterating gives the higher block
graphs \(X^{[n]}\).
\end{definition}

The target map \((u,a,v)\mapsto v\) is a right cover on the arcs. It is not a right cover in
the sense of Theorem~\ref{thm:rightcover}, because a vertex of \(X\) without incoming arcs has no
preimage. What replaces it is a shift.

\begin{lemma}\label{lem:shift}
\(\mathrm{ntot}_{X^{[2]}}(n+1)=\mathrm{ntot}_X(n+2)\). If a sequence \(u\) of naturals
satisfies \(u(n+1)\le c\,u(n)\), then \(u\) and \(n\mapsto u(n+1)\) have the same exponential
growth rate.
\end{lemma}

\begin{proof}
A path of length \(n+1\) from the arc \((u,a,v)\) is a path of length \(n+1\) from \(v\).
Summing over the arcs out of \(u\) gives \(\mathrm{npaths}_X(u,n+2)\). For the second claim,
\(\log u(n+1)/n\ge\log u(n+1)/(n+1)\) because \(u(n+1)\) is \(0\) or at least \(1\). The
\(\limsup\) of the right-hand side is that of \(\log u(n)/n\), and the bound \(u(n+1)\le c\,u(n)\)
gives the other inequality. Total path counts satisfy the bound with
\(c=|V|^2|A|\).
\end{proof}

\begin{theorem}\label{thm:linegraph}
The line graph has the entropy of \(X\), has zero entropy iff \(X\) has, and has an empty core
iff \(X\) has.
\end{theorem}

The window-\(K\) states of the \texttt{DubC2} engine and the contracted graphs of
\cite{Oda26} record several consecutive steps in one vertex. They are presentations of this
higher-block kind, and Theorem~\ref{thm:linegraph} is why recording more history never changes
a verdict.

\section{Homomorphisms and simulations}

A homomorphism \(f\colon X\to Y\) (Definition~\ref{def:carry-graph}) maps paths to paths with
the same labels. Three things follow at once.
\begin{itemize}
\item Cycles go to cycles, so a set containing every cycle of \(Y\) pulls back to one
  containing every cycle of \(X\). With a peel rank this is the rung of
  Corollary~\ref{cor:rung}.
\item Bi-infinite paths go to bi-infinite paths, so \(f\) maps the core into the core, and
  \(\operatorname{core}Y=\emptyset\) implies \(\operatorname{core}X=\emptyset\).
\item The label form passes from \(Y\) to \(X\) (Proposition~\ref{prop:state-label}(c)).
\end{itemize}
Path counts need one more condition.

\begin{proposition}\label{prop:locinj}\index{homomorphism!locally injective}
Let \(f\) be \emph{locally injective}: injective on the \(a\)-successors of each vertex. Then
\[
\mathrm{npaths}_X(v,n)\le\mathrm{npaths}_Y(f(v),n).
\]
So the entropy of \(X\) is at most
that of \(Y\), and zero entropy passes from \(Y\) to \(X\).
\end{proposition}

\begin{proof}
Induction on \(n\). Local injectivity turns the sum over the \(a\)-successors of \(v\) into a sum
over their images, which is a subset of the \(a\)-successors of \(f(v)\).
\end{proof}

A \emph{simulation}\index{simulation} of \(X\) by \(Y\) is a relation \(R\subseteq V\times W\) such that from
related vertices every arc \(v\xrightarrow{a}v'\) of \(X\) is matched by an arc
\(w\xrightarrow{a}w'\) of \(Y\) with \(R(v',w')\). A homomorphism is a functional simulation.
Infinite paths lift along a simulation one step at a time, so words of \(X\) are words of
\(Y\), and if every vertex of \(X\) is related to some vertex of \(Y\), the label form passes from
\(Y\) to \(X\).

\begin{theorem}[peel, then simulate]\label{thm:peel-simulate}
Let \(A_0,\dots,A_J\) be a peel chain of \(X\) (Definition~\ref{def:peel-chain}) with
\(A_J\subseteq K\). If the restriction of \(X\) to \(K\) is simulated by a graph \(Y\) with the
label form, and every vertex of \(K\) is related to a vertex of \(Y\), then \(X\) has the label
form.
\end{theorem}

\begin{proof}
An infinite path of \(X\) lies in \(K\) from step \(J\) on
(Proposition~\ref{prop:peel-chain}). Its tail lifts to \(Y\) with the same labels, which are
eventually periodic.
\end{proof}

This is the shape of a two-part certificate: rounds that push every infinite walk into a small
set \(K\), then a deterministic automaton on \(K\) with a rank.

\section{Label products}

\begin{definition}\label{def:labelproduct}
The \emph{label product}\index{label product} \(X\otimes Y\) \cite[Def.~3.4.8]{LM95} has vertex set \(V\times W\)
and an arc \((v,w)\xrightarrow{a}(v',w')\) whenever \(v\xrightarrow{a}v'\) and
\(w\xrightarrow{a}w'\). The indexed product \(\bigotimes_i X_i\) of a finite family is defined
in the same way.
\end{definition}

\begin{proposition}\label{prop:labelproduct}
\begin{enumerate}[label=(\alph*)]
\item A path of \(X\otimes Y\) is a pair of paths reading the same word. So
  \(L_{X\otimes Y}(v,w)=L_X(v)\cap L_Y(w)\), and sofic languages are closed under intersection
  \cite[Prop.~3.4.10]{LM95}. Products of right-resolving graphs are right-resolving.
\item The projections are homomorphisms. If one factor has the label form or an empty core,
  so does the product.
\item One step of the subset construction factorises, and likewise for indexed products:
  \[
  \operatorname{detStep}(S\times T,a)=\operatorname{detStep}(S,a)\times
  \operatorname{detStep}(T,a).
  \]
\item \(h(X\otimes Y)\le h(X)+h(Y)\). If \(Y\) is right-resolving, then \(h(X\otimes Y)\le h(X)\).
  If both factors have zero entropy, so does the product.
\item The number of words of length \(n\) from \((v,w)\) is at most the minimum of the
  numbers from \(v\) and from \(w\).
\end{enumerate}
\end{proposition}

\begin{proof}
(a), (b) and (c) are immediate from the definition. (d) By induction,
\[
\mathrm{npaths}((v,w),n)\le\mathrm{npaths}(v,n)\,\mathrm{npaths}(w,n),
\]
and the growth rate of a product of sequences is at most the sum of the rates. The entropy is finite
by the super-eigenvector bound of Theorem~\ref{thm:brackets} with the constant vector. If \(Y\)
is right-resolving, the first projection is locally injective (Proposition~\ref{prop:locinj}).
Polynomial bounds multiply. (e) follows from (a).
\end{proof}

The bound \(h(X\otimes Y)\le\min(h(X),h(Y))\) is false for path counts. Let \(X\) be a single
vertex with an \(a\)-loop, and \(Y\) two vertices with all four \(a\)-arcs. Then \(X\otimes Y\) is a
copy of \(Y\), with \(2^{n+1}\) paths of length \(n\) and entropy \(\log2>0=h(X)\). The bound
holds for words by (e), because words are what the labels see.

Part (c) is the graph content of the Chinese-remainder step in the follower automaton of
\texttt{DubC2}: the residue coordinate modulo \(M=\prod p_i^{k_i}\) is the indexed product of the
coordinates modulo \(p_i^{k_i}\), and the subset construction acts on each factor separately.
For an inclusion \(I\subseteq J\) of index sets, restriction is a homomorphism
\(\bigotimes_{j\in J}X_j\to\bigotimes_{i\in I}X_i\). So adding primes to \(\mathcal P\) shrinks
the language and keeps every certificate. The same holds next to a shared factor such as the
cell coordinate, because the product of two homomorphisms is a homomorphism of label products.

\section{Subdivision}\label{sec:subdivision}

A \emph{word graph} is a finite set \(E\subseteq V\times A^*\times V\) of arcs labelled by
non-empty words. An \emph{infinite word path} \((v_i,w_i)\) with times \(t_{i+1}=t_i+|w_i|\)
outputs the sequence \(\ell\) with \(\ell(t_i+k)=w_i[k]\). The graphs of \cite{Oda26} are word
graphs, and they become carry graphs by subdivision.

\begin{definition}\label{def:subdivide}
The \emph{subdivision}\index{subdivision} \(\operatorname{sd}E\) has the vertices of \(V\) and, for each arc
\(e=(u,w,v)\) and each \(0<i<|w|\), an interior vertex \((e,i)\). The arc \(e\) becomes the chain
\[
u=(e,0)\xrightarrow{w[0]}(e,1)\xrightarrow{w[1]}\cdots\xrightarrow{w[|w|-1]}(e,|w|)=v .
\]
\end{definition}

\begin{theorem}[subdivision]\label{thm:subdivide}
Let every word of \(E\) be non-empty.
\begin{enumerate}[label=(\roman*)]
\item Between vertices of \(V\), \(\operatorname{sd}E\) reads exactly the concatenations of the
  words along finite word paths.
\item \(\operatorname{sd}E\) has the label form iff every infinite word path of \(E\) has
  eventually periodic output.
\item \(\operatorname{sd}E\) has zero entropy iff no vertex of \(E\) has two distinct word arcs
  that both return to it.
\item The prefixes of the outputs of word paths from \(u\) form the language of
  \(\operatorname{sd}E\) at \(u\), so they form a sofic language.
\end{enumerate}
\end{theorem}

\begin{proof}
Every vertex reached from \(V\) is a position \((e,i)\) with \(0<i\le|w|\), and its rest
\(w[i..]\) leads to the end of \(e\). (i) follows by induction on the word read. (iii) An
interior vertex has a single out-arc, so branching can only happen at \(V\). Two distinct word
arcs out of \(u\) leave \(u\) along distinct pairs (label, target), and a target returns to \(u\)
iff the end of its word arc does. (iv) follows from (i). (ii) In one direction a word path
becomes a path of \(\operatorname{sd}E\) with the output as labels. In the other, an infinite
path of \(\operatorname{sd}E\) reaches \(V\) within \(\max|w|\) steps. From there it runs
through one whole word arc after another, and a recursive choice of these arcs gives a word
path with the same output.
\end{proof}

Part (ii) and the determinization theorem (Theorem~\ref{thm:determinization}) together extend
\cite[Thm.~21]{Oda26} from one-letter graphs (Chapter~\ref{ch:sofic}) to every word graph with
finitely many arcs. The phase test holds on
every strong component iff every infinite word path has eventually periodic output, iff
\(\det\operatorname{sd}E\) has the state form. This covers the contracted graphs of the normal
form. The one hypothesis is that the edge set is a finite set of non-empty words.

\section*{Notes}
State splitting, higher blocks and label products are \cite[§§2.3, 2.4, 3.4]{LM95}. Line
digraphs are \cite[§4.5]{BJG09}. Right covers are the graph side of the right-resolving factor
codes of \cite[ch.~8]{LM95}. The counting form of their invariance needs no conjugacy, no
Perron--Frobenius theory and no adjacency matrix. Example~\ref{ex:full-reduced} is gate X-SPLIT
of the certificate-graph study: the fibres of the full core over the reduced one have size
\(1\) or \(2\), namely the number of distinct incoming digits. Theorem~\ref{thm:subdivide}
answers gate G-2 of the plan positively, at the level of whole graphs. The per-component form,
in which the phase test on one strong component equals the state form of its subdivision, is
not needed.

\section*{Lean dictionary}
\begin{center}\small
\begin{tabular}{@{}>{\raggedright}p{0.3\textwidth}>{\raggedright\arraybackslash}p{0.64\textwidth}@{}}
\toprule
Text & Lean (\texttt{CarryGraph/Splitting.lean}, \texttt{Quotient.lean}, \texttt{Product.lean})\\
\midrule
Definition~\ref{def:rightcover} & \lean{RightCover}, \lean{RightCover.toHom}, \lean{IsSink}\\
Proposition~\ref{prop:rightcover-paths} & \lean{npaths\_eq\_of\_rightCover}\\
Theorem~\ref{thm:rightcover} & \lean{ntot\_le\_of\_rightCover}, \lean{ntot\_rightCover\_le},
  \lean{entropy\_eq\_of\_ntot}, \lean{core\_eq\_empty\_iff\_of\_ntot},
  \lean{entropy\_eq\_of\_rightCover}, \lean{lnoBranch\_iff\_of\_rightCover},
  \lean{polyGrowth\_iff\_of\_rightCover}, \lean{core\_eq\_empty\_iff\_of\_rightCover}\\
Example~\ref{ex:full-reduced} & \lean{CarryBridge.reducedGraph},
  \lean{CarryBridge.digitGraph\_rightCover}, \lean{CarryBridge.npaths\_full\_eq\_reduced},
  \lean{CarryBridge.exists\_unit\_lift}, \lean{CarryBridge.certOK\_full\_iff\_reduced'}\\
Definition~\ref{def:split}, Theorem~\ref{thm:split} & \lean{inSplit}, \lean{adj\_inSplit},
  \lean{inSplit\_rightCover}, \lean{npaths\_inSplit}, \lean{entropy\_inSplit},
  \lean{lnoBranch\_inSplit\_iff}, \lean{core\_inSplit}; \lean{outSplit}, \lean{adj\_outSplit},
  \lean{sum\_npaths\_outSplit}, \lean{ntot\_outSplit}, \lean{entropy\_outSplit},
  \lean{lnoBranch\_outSplit\_iff}, \lean{core\_outSplit}\\
Definition~\ref{def:linegraph} & \lean{lineGraph}, \lean{adj\_lineGraph},
  \lean{lineGraph\_rightCover}, \lean{npaths\_lineGraph}\\
Lemma~\ref{lem:shift} & \lean{ntot\_lineGraph}, \lean{expGrowthSup\_shift},
  \lean{ntot\_succ\_le}, \lean{core\_eq\_empty\_iff\_ntot}\\
Theorem~\ref{thm:linegraph} & \lean{entropy\_lineGraph}, \lean{lnoBranch\_lineGraph\_iff},
  \lean{core\_lineGraph}\\
Homomorphisms & \lean{Hom.mem\_core}, \lean{core\_subset\_preimage\_core},
  \lean{Hom.core\_eq\_empty}, \lean{cyc\_mem\_of\_peelRank\_of\_hom}\\
Proposition~\ref{prop:locinj} & \lean{Hom.LocInj}, \lean{Hom.npaths\_le},
  \lean{Hom.polyGrowth\_of\_locInj}, \lean{Hom.entropy\_le\_of\_locInj},
  \lean{Hom.lnoBranch\_of\_locInj}\\
Simulations & \lean{IsSimulation}, \lean{Hom.isSimulation}, \lean{IsSimulation.reads},
  \lean{IsSimulation.langFrom\_subset}, \lean{IsSimulation.exists\_iPath},
  \lean{allPer\_of\_simulation}\\
Theorem~\ref{thm:peel-simulate} & \lean{allPer\_of\_peelChain}\\
\bottomrule
\end{tabular}
\end{center}

\begin{center}\small
\begin{tabular}{@{}>{\raggedright}p{0.3\textwidth}>{\raggedright\arraybackslash}p{0.64\textwidth}@{}}
\toprule
Text & Lean (\texttt{CarryGraph/Product.lean}, \texttt{Quotient.lean})\\
\midrule
Definition~\ref{def:labelproduct} & \lean{labelProduct}, \lean{adj\_labelProduct},
  \lean{fstHom}, \lean{sndHom}, \lean{Hom.prodMap}; \lean{piProduct}, \lean{adj\_piProduct},
  \lean{projHom}\\
Proposition~\ref{prop:labelproduct} & (a) \lean{reads\_labelProduct},
  \lean{mem\_langFrom\_labelProduct}, \lean{langFrom\_labelProduct}, \lean{Sofic.inter},
  \lean{rightResolving\_labelProduct}, \lean{reads\_piProduct},
  \lean{mem\_langFrom\_piProduct}, \lean{rightResolving\_piProduct}; (b)
  \lean{allPer\_labelProduct\_left}, \lean{allPer\_labelProduct\_right},
  \lean{core\_labelProduct\_left}, \lean{core\_labelProduct\_right}, \lean{allPer\_piProduct};
  (c) \lean{detStep\_labelProduct}, \lean{detStep\_piProduct}; (d)
  \lean{npaths\_labelProduct\_le}, \lean{ntot\_labelProduct\_le}, \lean{entropy\_ne\_top},
  \lean{entropy\_labelProduct\_le}, \lean{fstHom\_locInj}, \lean{entropy\_labelProduct\_le\_left},
  \lean{lnoBranch\_labelProduct}; (e) \lean{card\_words\_labelProduct\_le}\\
Monotonicity & \lean{restrictHom}, \lean{allPer\_piProduct\_mono},
  \lean{langFrom\_piProduct\_antitone}, \lean{allPer\_labelProduct\_piProduct\_mono}\\
Word graphs & \lean{WReads}, \lean{WReach}, \lean{WNoBranch}, \lean{WCovers}, \lean{WAllPer}\\
Definition~\ref{def:subdivide} & \lean{sdLen}, \lean{sdPos}, \lean{subdivide},
  \lean{adj\_subdivide}, \lean{sd\_step\_inl}, \lean{sd\_step\_inr}, \lean{reads\_sdPos}\\
Theorem~\ref{thm:subdivide} & (i) \lean{reads\_subdivide\_iff}, \lean{wReads\_of\_reads};
  (ii) \lean{allPer\_subdivide\_iff}, \lean{sd\_follow}, \lean{sd\_arc}; (iii)
  \lean{lnoBranch\_subdivide\_iff}, \lean{eq\_of\_sdPos\_one}; (iv)
  \lean{mem\_langFrom\_subdivide}, \lean{sofic\_wordLang}\\
\cite[Thm.~21]{Oda26}, all finite word graphs & \lean{wAllPer\_iff\_det},
  \lean{CarryBridge.covers\_iff\_wCovers}, \lean{CarryBridge.odaAllPer\_iff\_subdivide},
  \lean{CarryBridge.allCert\_iff\_subdivide}\\
\bottomrule
\end{tabular}
\end{center}
